\section{Decisiones pendientes}

\subsection{Descripcion}

Este documento concentra las decisiones aun abiertas del proyecto en un registro unico, visible y trazable. Su funcion es evitar que reglas importantes queden dispersas en conversaciones o notas informales.

\subsection{Alcance}

Aplica a reglas de negocio, arquitectura, datos, integraciones y criterios operativos que todavia no han sido cerrados formalmente.

\subsection{Definiciones}

\begin{itemize}
  \item \textbf{Decision pendiente}: definicion necesaria aun no cerrada formalmente.
  \item \textbf{Bloqueante}: decision cuya ausencia impide documentar o construir correctamente una parte relevante del sistema.
  \item \textbf{Riesgo}: decision abierta que no bloquea por completo, pero puede introducir retrabajo, inconsistencia o deuda de producto.
  \item \textbf{Criticidad}: alcance del impacto de una decision abierta sobre frontend, backend o ambos.
\end{itemize}

\subsection{Reglas}

\begin{itemize}
  \item Toda decision pendiente debe clasificarse por tipo y criticidad.
  \item Una decision cerrada debe actualizar su documento fuente y el changelog.
  \item La version mas reciente aprobada en este repositorio prevalece sobre notas previas y conversaciones no integradas al registro.
  \item Este registro no reemplaza al documento fuente de cada regla; solo centraliza el estado de las definiciones abiertas.
\end{itemize}

\subsection{Registro activo}

\subsubsection{Prioridad alta}

\begin{enumerate}
  \item \textbf{Consolidacion del flujo de comprobantes con Google Forms}\\
  Tipo: \texttt{Decision pendiente}\\
  Criticidad: \texttt{Bloquea activacion y release}\\
  Documento responsable: \texttt{10\_01\_Roadmap\_y\_Fases}, \texttt{01\_02\_Estados\_y\_Flujos} y \texttt{03\_02\_Procesos\_de\_Negocio}\\
  Motivo: ya se decidio abandonar upload de comprobantes en la plataforma, pero falta cerrar el contrato operativo entre registro en la app, evidencia enviada por Google Forms y aprobacion/rechazo administrativa en backend.
\end{enumerate}

\subsubsection{Prioridad media}

\begin{enumerate}
  \setcounter{enumi}{1}
  \item \textbf{Alcance final del track de migracion a Supabase Auth}\\
  Tipo: \texttt{Riesgo}\\
  Criticidad: \texttt{No bloqueante para roadmap competitivo}\\
  Documento responsable: \texttt{10\_01\_Roadmap\_y\_Fases}\\
  Motivo: la migracion es trabajo real en curso, pero no tiene sprint competitivo fijo. Falta decidir solo su cierre operativo final, no su pertenencia al goal de los sprints competitivos.
\end{enumerate}

\subsubsection{Prioridad baja}

\begin{enumerate}
  \setcounter{enumi}{2}
  \item \textbf{Persistencia explicita de \texttt{SIN\_PICK}}\\
  Tipo: \texttt{Decision pendiente}\\
  Criticidad: \texttt{No bloqueante}\\
  Documento responsable: \texttt{01\_03\_Picks\_y\_Locks} y \texttt{05\_Datos}\\
  Motivo: la regla funcional ya esta cerrada, pero falta decidir si se persiste como estado explicito o si solo se deriva del negocio.

  \item \textbf{Distincion visual entre correccion por API y override manual}\\
  Tipo: \texttt{Riesgo}\\
  Criticidad: \texttt{No bloqueante}\\
  Documento responsable: \texttt{01\_05\_Casos\_Borde}\\
  Motivo: la prioridad operativa ya esta definida, pero falta decidir si la interfaz hace visible el origen de la correccion.

  \item \textbf{Version resumida del reglamento para participantes}\\
  Tipo: \texttt{Riesgo}\\
  Criticidad: \texttt{No bloqueante}\\
  Documento responsable: \texttt{01\_01\_Reglamento\_V1}\\
  Motivo: el reglamento actual sirve como fuente normativa, pero aun no se decide si se publicara una version mas breve para consumo general.
\end{enumerate}

\subsection{Decisiones cerradas recientemente}

\begin{itemize}
  \item \textbf{Superficie del leaderboard en Sprint 6}\\
  Estado: \texttt{Cerrada en Sprint 6}\\
  Veredicto: una sola superficie autenticada de ranking con filtros para general, fase y jornada.\\
  Registro de cierre: \texttt{12\_Decisiones\_por\_Sprint/12\_06\_Sprint\_06.tex}

  \item \textbf{Precision visual y desglose de scoring en Sprint 6}\\
  Estado: \texttt{Cerrada en Sprint 6}\\
  Veredicto: enteros sin decimales, fraccionarios con dos decimales visibles; el total es el dato primario y el desglose es contexto secundario.\\
  Registro de cierre: \texttt{12\_Decisiones\_por\_Sprint/12\_06\_Sprint\_06.tex} y \texttt{01\_Producto\_y\_Reglas/tex/01\_04\_Scoring\_y\_Desempates.tex}

  \item \textbf{Tratamiento de picks especiales en leaderboard filtrado}\\
  Estado: \texttt{Cerrada en Sprint 6}\\
  Veredicto: la vista filtrada por jornada o fase excluye picks especiales del puntaje principal visible.\\
  Registro de cierre: \texttt{12\_Decisiones\_por\_Sprint/12\_06\_Sprint\_06.tex}

  \item \textbf{Feedback visible de resultado corregido en Sprint 6}\\
  Estado: \texttt{Cerrada en Sprint 6}\\
  Veredicto: badge persistente en el partido corregido y aviso breve cuando cambie puntaje o ranking.\\
  Registro de cierre: \texttt{12\_Decisiones\_por\_Sprint/12\_06\_Sprint\_06.tex}

  \item \textbf{Modelo de recalculo minimo de Sprint 6}\\
  Estado: \texttt{Cerrada en Sprint 6}\\
  Veredicto: recalculo automatico selectivo por partido afectado; el recalculo global admin queda fuera del cierre minimo del sprint.\\
  Registro de cierre: \texttt{12\_Decisiones\_por\_Sprint/12\_06\_Sprint\_06.tex} y \texttt{03\_Backend/tex/03\_02\_Procesos\_de\_Negocio.tex}

  \item \textbf{Ubicacion temporal del refactor de comprobantes con Google Forms}\\
  Estado: \texttt{Cerrada}\\
  Veredicto: el refactor del flujo de comprobantes para abandonar upload en plataforma y consolidar evidencia via Google Forms se planifica para Sprint 7.\\
  Registro de cierre: \texttt{10\_Planeacion/tex/10\_01\_Roadmap\_y\_Fases.tex}

  \item \textbf{Tratamiento del track de migracion a Supabase Auth}\\
  Estado: \texttt{Cerrada}\\
  Veredicto: se reconoce como track paralelo real sin sprint competitivo fijo y no redefine por si solo el goal de los sprints competitivos.\\
  Registro de cierre: \texttt{10\_Planeacion/tex/10\_01\_Roadmap\_y\_Fases.tex}

  \item \textbf{Presentacion del flujo separado de predicciones especiales}\\
  Estado: \texttt{Cerrada en Sprint 3}\\
  Veredicto: las predicciones especiales viven en una vista separada dentro del modulo funcional de picks y no se mezclan con la captura de picks por partido.\\
  Registro de cierre: \texttt{12\_Decisiones\_por\_Sprint/12\_03\_Sprint\_03.pdf}

  \item \textbf{Politica de visibilidad de Campeon y Goleador frente a otros usuarios}\\
  Estado: \texttt{Cerrada en Sprint 3}\\
  Veredicto: permanecen ocultos para otros usuarios hasta que exista una regla oficial posterior de revelacion.\\
  Registro de cierre: \texttt{12\_Decisiones\_por\_Sprint/12\_03\_Sprint\_03.pdf}

  \item \textbf{Contrato publico de consulta de resultado oficial de partido (QH-107)}\\
  Estado: \texttt{Cerrada en Sprint 5}\\
  Veredicto: la consulta publica acordada es \texttt{GET /matches/\{matchId\}/result}; requiere autenticacion; responde \texttt{200 OK} con \texttt{result = null} si aun no existe resultado oficial; y expone solo \texttt{matchId}, \texttt{result.homeGoals}, \texttt{result.awayGoals} y \texttt{result.updatedAt}. Los campos internos de trazabilidad y fuente permanecen fuera del contrato publico.\\
  Registro de cierre: \texttt{03\_Backend/tex/03\_04\_Endpoints\_Conceptuales.tex} y \texttt{03\_Backend/tex/03\_02\_Procesos\_de\_Negocio.tex}
\end{itemize}

\subsection{Edge Cases}

\begin{itemize}
  \item Una decision puede bajar de criticidad si el documento fuente ya define un comportamiento por defecto suficiente para desarrollo interno.
  \item Una decision cerrada en conversacion no se considera resuelta hasta que quede reflejada en su documento fuente y en este registro.
\end{itemize}

\subsection{Invariantes}

\begin{itemize}
  \item No debe haber decisiones relevantes abiertas solo en conversaciones informales.
  \item Toda decision abierta de producto debe tener un documento responsable identificable.
  \item La lista priorizada debe poder leerse sin necesidad de revisar todos los documentos fuente.
\end{itemize}

\subsection{Preguntas abiertas}

\begin{itemize}
  \item No hay preguntas abiertas sobre el formato del registro en esta fase; el siguiente paso es mantenerlo sincronizado conforme avancemos documento por documento.
\end{itemize}
